
\section{Method}

\subsection{Overview}

AASVD compresses each targeted \texttt{nn.Linear} (or \texttt{Conv1D}) layer in four phases:

\begin{enumerate}
\item \textbf{Calibration} ( 4.2): collect $C_X$ per layer via streaming EMA.
\item \textbf{Weighted SVD} ( 4.3): factorize $W C_X^{1/2}$, then unweight.
\item \textbf{Validation} ( 4.4): compute Frobenius reconstruction error.
\item \textbf{Replacement} ( 4.5): substitute \texttt{SVDLinear} if error is within threshold.

\end{enumerate}
\textbf{Figure 4.1} shows the pipeline.

\begin{lstlisting}
   +-------------------------------------------+
   |  Phase 1: Streaming EMA Calibration       |
   |  Input: bundled or user texts             |
   |  Output: C\_X^( ) for each target layer    |
   +----------------+--------------------------+
                    |
                     
   +-------------------------------------------+
   |  For each layer  :                        |
   |    Phase 2: Weighted SVD at rank r        |
   |    Phase 3: Frobenius validation          |
   |    Phase 4: Replace or fallback           |
   +----------------+--------------------------+
                    |
                     
   +-------------------------------------------+
   |  Compressed model with SVDLinear layers   |
   |  Ready for LoRA recovery fine-tuning      |
   +-------------------------------------------+
\end{lstlisting}

\subsection{Phase 1: Streaming Covariance Calibration}

\subsubsection{Target Statistic}

For each target linear layer $\ell$, we want:

\[
C\_X^{(\ell)} = \mathbb{E}\_{x^{(\ell)} \sim \mathcal{D}}\!\left[x^{(\ell)} (x^{(\ell)})^\top\right]
\tag{4.1}
\]

where $x^{(\ell)}$ is the layer input during calibration.

\subsubsection{EMA Update}

Given a batch of $B_\text{calib}$ calibration inputs $X \in \mathbb{R}^{B_\text{calib} \times d}$, 
the batch covariance is:

\[
C\_\text{batch} = \frac{1}{N\_\text{valid}} \sum\_{i=1}^{N\_\text{valid}} x\_i x\_i^\top
\tag{4.2}
\]

where $N_\text{valid} \leq B_\text{calib}$ counts non-padding tokens. We 
then update the running estimate:

\[
\hat{C}\_X \leftarrow \beta \hat{C}\_X + (1 - \beta) C\_\text{batch}
\tag{4.3}
\]

with $\beta = 0.9$ by default. The running estimate starts as the first 
batch covariance $\hat{C}_X^{(0)} = C_\text{batch}^{(0)}$.

\subsubsection{Padding Token Filtering}

To avoid bias from padding tokens (which are typically zero):

\[
\text{valid}(x\_i) = \begin{cases} 1 & \text{if } \|x\_i\|\_1 \textgreater{} 10^{-5} \\ 0 & \text{otherwise} \end{cases}
\tag{4.4}
\]

Only valid tokens contribute to $C_\text{batch}$. This is critical for 
batched calibration, where padding dominates the raw token count.

\subsubsection{Memory Layout}

We compute $C_\text{batch}$ on GPU (fast matmul), then **move it to CPU 
immediately**:

\begin{lstlisting}
batch\_cov = (torch.matmul(x\_active.T, x\_active) / max(1, x\_active.size(0))).cpu()
\end{lstlisting}

This prevents GPU memory accumulation over calibration iterations. The total 
GPU footprint of calibration is one batch's activation, plus the running 
covariance estimate (kept on CPU).

\subsubsection{Calibration Corpus}

We ship a \textbf{default multi-domain corpus} of 5 documents:

\begin{table}[H]
\centering
\small
\begin{tabular}{lll}
\toprule
\textbf{\#} & \textbf{Domain} & \textbf{Approximate length} \\
\midrule
1 \& Technical prose (attention mechanisms) & 45 tokens \\
2 \& Mathematical prose (SVD, activations) & 38 tokens \\
3 \& Source code (Python attention) & 42 tokens \\
4 \& Technical prose (distributed training) & 42 tokens \\
5 \& Mathematical prose (chain-of-thought) & 48 tokens \\
\bottomrule
\end{tabular}
\end{table}

Total: ~215 tokens. The corpus is designed to activate diverse subspaces 
of the residual stream, providing a broad calibration signal even without 
access to the target task.

If user-supplied \texttt{calibration\_texts} are provided, they \textbf{replace} the 
default corpus.

\subsection{Phase 2: Weighted SVD}

\subsubsection{The Weighted Objective}

Standard SVD minimizes $\|W - W_r\|_F^2$, treating all input dimensions 
equally. We instead minimize the activation-weighted error:

\[
\mathcal{L}\_\text{weighted}(W\_r) = \mathbb{E}\_{x}\!\left[\|W x - W\_r x\|\_2^2\right]
\tag{4.5}
\]

Under the second-moment assumption $\mathbb{E}[x x^\top] = C_X$:

\[
\mathcal{L}\_\text{weighted}(W\_r) = \mathrm{Tr}\!\left[(W - W\_r) C\_X (W - W\_r)^\top\right] = \|(W - W\_r) C\_X^{1/2}\|\_F^2
\tag{4.6}
\]

which is exactly an unweighted problem on the transformed matrix:

\[
W\_\text{weighted} = W C\_X^{1/2}
\tag{4.7}
\]

\subsubsection{The Algorithm}

\begin{enumerate}
\item Compute $W_\text{weighted} = W C_X^{1/2}$.
\item Compute the truncated SVD:
  $$
  W\_\text{weighted} \approx U\_r \Sigma\_r V\_r^\top
  \tag{4.8}
  $$
\item Unweight the right factor:
  $$
  V\_r^\text{unweighted} = V\_r^\top C\_X^{-1/2}
  \tag{4.9}
  $$
\item Distribute the singular values symmetrically:
  $$
  U = U\_r \Sigma\_r^{1/2}, \qquad
  V = \Sigma\_r^{1/2} V\_r^\text{unweighted}
  \tag{4.10}
  $$
\item The factorized layer computes:
  $$
  y = (x V^\top) U^\top + b
  \tag{4.11}
  $$

\end{enumerate}
\subsubsection{Eigendecomposition with Regularization}

We compute $C_X^{1/2}$ and $C_X^{-1/2}$ via eigendecomposition:

\[
C\_X = Q \Lambda Q^\top, \quad
C\_X^{1/2} = Q \Lambda^{1/2} Q^\top, \quad
C\_X^{-1/2} = Q \Lambda^{-1/2} Q^\top
\tag{4.12}
\]

Eigenvalues are regularized:

\[
\lambda\_i \leftarrow \max(\lambda\_i, \epsilon\_\text{reg}), \qquad
\epsilon\_\text{reg} = 10^{-6} \cdot \max\_j \lambda\_j + 10^{-8}
\tag{4.13}
\]

This ensures $C_X^{-1/2}$ is finite for near-zero eigenvalues.

\subsubsection{Rank Selection}

We use a rank constrained by \textbf{both} the target rank $r$ and a parameter 
budget constraint:

\[
r\_\text{eff} = \min\!\left(r, \left\lfloor 0.75 \cdot \frac{d\_\text{in} d\_\text{out}}{d\_\text{in} + d\_\text{out}} \right\rfloor \right)
\tag{4.14}
\]

The $0.75$ factor enforces that the factorized layer uses at most 75\% of the 
original parameters, guaranteeing meaningful compression.

If $r_\text{eff} < r$, the requested rank is reduced; this is common for 
narrow layers (e.g., $d_\text{in} = d_\text{out} = 128$ at $r = 128$).

\subsection{Phase 3: Frobenius Validation}

After factorization, we compute:

\[
\text{err}(W, U, V) = \frac{\|W - UV\|\_F}{\|W\|\_F + 10^{-12}}
\tag{4.15}
\]

\textbf{Two possible outcomes:}

\begin{enumerate}
\item \textbf{err $\leq \tau$.} The factorization is accepted. \texttt{SVDLinear} replaces
  the original layer.
\item \textbf{err $> \tau$.} If \texttt{fallback\_to\_original = True}, the original layer
  is preserved and logged. Otherwise, the factorization is applied anyway
  with a warning.

\end{enumerate}
We default to \texttt{fallback\_to\_original = True} and \texttt{max\_error = 0.50}.

\subsection{Phase 4: Structural Replacement}

\subsubsection{The `SVDLinear` Module}

The replacement layer implements:

\[
y = (x V^\top) U^\top + b + \text{LoRA}(x) \cdot s
\tag{4.16}
\]

where:

\begin{itemize}
\item $V \in \mathbb{R}^{r \times d_\text{in}}$: down-projection
\item $U \in \mathbb{R}^{d_\text{out} \times r}$: up-projection
\item $b \in \mathbb{R}^{d_\text{out}}$: bias (if original had one)
\item $\text{LoRA}(x) = (x A^\top) B^\top$: optional recovery adapter

\end{itemize}
\subsubsection{Parameter Count}

\begin{table}[H]
\centering
\small
\begin{tabular}{ll}
\toprule
\textbf{Configuration} & \textbf{Parameters} \\
\midrule
Original & $d_\text{in} d_\text{out}$ \\
AASVD at rank $r$ & $r(d_\text{in} + d_\text{out})$ \\
AASVD + LoRA($r_L$) & $r(d_\text{in} + d_\text{out}) + r_L(d_\text{in} + d_\text{out})$ \\
\bottomrule
\end{tabular}
\end{table}

For a $768 \times 768$ matrix at $r = 64$:

\begin{table}[H]
\centering
\small
\begin{tabular}{lll}
\toprule
\textbf{Config} & \textbf{Params} & \textbf{Ratio} \\
\midrule
Original & 589,824 & 1.00$\times$ \\
AASVD r=64 & 98,304 & \textbf{6.0$\times$} \\
AASVD r=64 + LoRA r=16 & 122,880 & \textbf{4.8$\times$} \\
\bottomrule
\end{tabular}
\end{table}

\subsubsection{LoRA-Freeze Semantics}

When LoRA is enabled, the base factors $U, V$ are \textbf{frozen} and only LoRA 
adapters $A, B$ are trainable. This ensures that:

\begin{enumerate}
\item The compression is not undone during recovery.
\item The parameter reduction persists after fine-tuning.
\item Gradient flow through $U, V$ is blocked to avoid drift.

\end{enumerate}
The freeze is implemented via \texttt{requires\_grad = False} on $U, V$.

\subsection{The Full Algorithm}

\textbf{Algorithm 1:} AASVD compression
\begin{lstlisting}
Input:  model M, rank r, threshold $\tau$, $\beta$=0.9
        calibration texts D\_calib
Output: compressed model M', statistics

1:  for   in target\_layers(M) do
2:      register\_forward\_hook( , accumulate\_C\_X)
3:  end for
4:
5:  # Phase 1: Calibration
6:  inputs $\leftarrow$ tokenize(D\_calib)
7:  forward(M, inputs)
8:  remove\_hooks()
9:
10: # Cache eigendecomposition per layer
11: for   in target\_layers(M) do
12:     Q,   $\leftarrow$ eigh(C\_X^( ))
13:     C\_sqrt[ ] $\leftarrow$ Q @ diag($\sqrt{}$ ) @ Q.T
14:     C\_inv\_sqrt[ ] $\leftarrow$ Q @ diag(1/$\sqrt{}$ ) @ Q.T
15: end for
16:
17: # Phase 2-4: Weighted SVD + validation
18: for   in target\_layers(M) do
19:     W $\leftarrow$ weight( )
20:     W\_weighted $\leftarrow$ W @ C\_sqrt[ ]
21:     U,  , V $\leftarrow$ svd(W\_weighted)
22:     U\_r, V\_r $\leftarrow$ truncate(U,  , V, r)
23:     V\_r $\leftarrow$ V\_r @ C\_inv\_sqrt[ ]
24:     W\_rec $\leftarrow$ U\_r @ V\_r
25:     err $\leftarrow$  W - W\_rec \_F / ( W \_F + $\epsilon$)
26:
27:     if err $\le$ $\tau$:
28:         replace( , SVDLinear(U\_r, V\_r))
29:     elif fallback\_to\_original:
30:         log\_warning( , err)
31:         mark\_fallback( )
32:     else:
33:         log\_warning( , err)
34:         replace( , SVDLinear(U\_r, V\_r))
35: end for
36:
37: return M'
\end{lstlisting}

\subsection{Complexity Analysis}

\begin{table}[H]
\centering
\small
\begin{tabular}{lll}
\toprule
\textbf{Phase} & \textbf{Time} & \textbf{GPU Memory} \\
\midrule
Calibration & $\mathcal{O}(N B d^2)$ & $\mathcal{O}(B d)$ \\
Eigendecomposition & $\mathcal{O}(N_l d^3)$ & $\mathcal{O}(d^2)$ \\
Weighted SVD & $\mathcal{O}(N_l d^3)$ & $\mathcal{O}(d^2)$ \\
Validation & $\mathcal{O}(N_l r d^2)$ & $\mathcal{O}(d^2)$ \\
\textbf{Total} & $\mathcal{O}(N_l d^3)$ & $\mathcal{O}(N_l d^2)$ \\
\bottomrule
\end{tabular}
\end{table}

CPU memory for storing all covariances: $\mathcal{O}(N_l d^2)$. For 
DistilBERT ($N_l = 6$, $d = 768$): ~14M values, ~56 MB at FP32.

\subsection{Theoretical Notes}

\subsubsection{Optimality of Weighted SVD}

\textbf{Proposition.} $W C_X^{1/2}$ truncated at rank $r$ minimizes $\|(W - W_r) C_X^{1/2}\|_F$ 
over all rank-$r$ matrices $W_r$.

\textbf{Proof.} This is a direct application of Eckart--Young to the matrix 
$W C_X^{1/2}$. $\square$

\subsubsection{Relation to Standard SVD}

If $C_X = I$ (identity, uncorrelated unit-variance inputs), then 
$W C_X^{1/2} = W$ and AASVD reduces to standard SVD.

If $C_X = \mathrm{diag}(F)$ (independent but non-uniform variance), then 
$W C_X^{1/2} = W \mathrm{diag}(\sqrt{F})$ and AASVD reduces to A3's 
diagonal weighting.

\textbf{AASVD is strictly more general} than both standard SVD and diagonal 
weighting.

\subsubsection{Error Bound}

\textbf{Proposition.} For AASVD with rank $r$:

\[
\|W - U\_r V\_r\|\_F \leq \frac{\sqrt{\sum\_{i \textgreater{} r} \sigma\_i^2}}{\sqrt{\lambda\_\text{min}(C\_X)}}
\tag{4.17}
\]

where $\sigma_i$ are the singular values of $W C_X^{1/2}$ and $\lambda_\text{min}$ 
is the smallest eigenvalue of $C_X$ (before regularization).

\textbf{Proof.} We have:

\[
\|W - U\_r V\_r\|\_F \cdot \sqrt{\lambda\_\text{min}(C\_X)}
\leq \|(W - U\_r V\_r) C\_X^{1/2}\|\_F
\leq \sqrt{\sum\_{i \textgreater{} r} \sigma\_i^2}
\tag{4.18}
\]

Dividing gives (4.17). $\square$

\subsubsection{EMA Convergence}

\textbf{Proposition.} The EMA estimate $\hat{C}_X^{(k)}$ converges to the true 
$\mathbb{E}[C_\text{batch}]$ as $k \to \infty$ with rate $\mathcal{O}(\beta^k)$.

\textbf{Proof sketch.} The EMA recursion $\hat{C}^{(k)} = \beta \hat{C}^{(k-1)} + (1-\beta) C_k$ 
has the closed form:

\[
\hat{C}^{(k)} = \beta^{k-1} C\_1 + (1 - \beta) \sum\_{j=2}^{k} \beta^{k-j} C\_j
\tag{4.19}
\]

Taking expectations (assuming $\mathbb{E}[C_k] = \bar{C}$):

\[
\mathbb{E}[\hat{C}^{(k)}] = \bar{C} + \beta^{k-1} (C\_1 - \bar{C})
\tag{4.20}
\]

so the bias decays as $\beta^{k-1}$. $\square$